\section*{Bab \ref{ch:b2:reaction-free-energy} --- Entropi dan Energi Bebas Reaksi}

\begin{solution}{exo:b2:reaction-free-energy:1}
(a) Tiga mol gas menjadi cairan: $\Delta_r S^\circ < 0$. (b) Terbentuk gas:
$> 0$. (c) Satu molekul gas menghasilkan dua: $> 0$. (d) Ion-ion dalam
larutan membentuk kristal: $< 0$. (e) Peleburan: $> 0$. (f)
$\Delta\nu_{\text{gas}} = 0$: kecil, dan di sini positif
(\qty{+20}{J/(K.mol)} menurut tabel, karena \ce{HCl} kurang simetris
daripada \ce{H2} dan \ce{Cl2}).
% ledger: s0:HCl_g, s0:H2_g, s0:Cl2_g
\end{solution}

\begin{solution}{exo:b2:reaction-free-energy:2}
$\Delta_r S^\circ = 2(192.77) - 191.61 - 3(130.68) = \qty{-198.1}{J/(K.mol)}$;
$\Delta_r G^\circ = -91.9 - 298.15 \times (-0.1981) = \qty{-32.8}{kJ/mol}$:
\omterm{def:b2:reaction-free-energy:exergonic}{eksergonik} pada \qty{298}{K}, walaupun suku entropi (\qty{+59.1}{kJ/mol})
melawannya.
\end{solution}

\begin{solution}{exo:b2:reaction-free-energy:3}
$\Delta_r S^\circ = 69.95 - 130.68 - \tfrac12(205.15) = \qty{-163.3}{J/(K.mol)}$;
$\Delta_r G^\circ = -285.83 + 298.15 \times 0.1633 = \qty{-237.14}{kJ/mol}$.
Reaksi ini adalah reaksi pembentukan air cair (satu mol dari unsur-unsur
dalam keadaan acuannya), sehingga \omterm{def:b2:reaction-free-energy:reaction-gibbs}{energi Gibbs reaksi standarnya} menurut
definisi adalah $\Delta_f G^\circ(\ce{H2O}, \text{l})$; tabel memberikan
nilai yang sama.
\end{solution}

\begin{solution}{exo:b2:reaction-free-energy:4}
Sekitar \qty{42}{J/(K.mol)} pada \qty{298}{K} dan \qty{87}{J/(K.mol)} pada
\qty{1000}{K} (cair). Lonjakan: \qty{10.6}{J/(K.mol)} pada \qty{692.7}{K}
dan \qty{97.7}{J/(K.mol)} pada \qty{1180.2}{K}. $\Delta_{\text{fus}}H = 10.6 \times
692.7 = \qty{7.3}{kJ/mol}$ dan $\Delta_{\text{vap}}H = 97.7 \times 1180.2 =
\qty{115}{kJ/mol}$.
\end{solution}

\begin{solution}{exo:b2:reaction-free-energy:5}
$\Delta_r H^\circ = \qty{44.00}{kJ/mol}$, $\Delta_r S^\circ = 188.84 - 69.95 =
\qty{118.9}{J/(K.mol)}$, $\Delta_r G^\circ = 44.00 - 298.15 \times 0.1189 =
\qty{+8.55}{kJ/mol}$. Pada \qty{298}{K} perubahan itu tidak reversibel (air
cair dan uap pada \qty{1}{bar} tidak berada dalam kesetimbangan), sehingga
$\Delta S \ne
Q/T$: $\Delta_r H^\circ/T = \qty{147.6}{J/(K.mol)}$.
$\Delta_r G^\circ$ yang positif berarti bahwa pada \qty{298}{K} air tidak
menguap ke dalam atmosfer uap murni pada \qty{1}{bar}: uap itulah yang
mengembun. Keduanya berada dalam kesetimbangan ketika $\Delta_r G^\circ = 0$:
$T = 44\,004/118.9 =
\qty{370}{K}$, hanya \qty{3}{K} dari titik didih
normal (\qty{373}{K}): \omterm{def:b2:reaction-free-energy:ellingham-approximation}{pendekatan Ellingham} baik dalam rentang
\qty{75}{K}.
\end{solution}

\begin{solution}{exo:b2:reaction-free-energy:6}
$T_i = 57\,100/175.7 = \qty{325}{K}$. Di bawahnya $\Delta_r G^\circ > 0$
dan \ce{N2O4} yang tidak berwarna lebih disukai; pendinginan tabung
menggeser gas ke arah \ce{N2O4} dan warna cokelat \ce{NO2} memudar
(hubungan kuantitatif dengan komposisi ada di
\cref{ch:b2:equilibrium-shifts}).
\end{solution}

\begin{solution}{exo:b2:reaction-free-energy:7}
$\Delta_r S^\circ = -\dd\Delta_r G^\circ/\dd T = -(-4.0 + 12.0)/100 =
\qty{-0.080}{kJ/(K.mol)} = \qty{-80}{J/(K.mol)}$ dan $\Delta_r H^\circ =
\Delta_r G^\circ + T\Delta_r S^\circ = -12.0 + 300 \times (-0.080) =
\qty{-36.0}{kJ/mol}$. Pemeriksaan: $\Delta_r G^\circ/T$ berubah dari
$-0.0400$ menjadi $\qty{-0.0100}{kJ/(K.mol)}$, kemiringan
$3.00 \times 10^{-4}$ per kelvin, sama dengan
$-\Delta_r H^\circ/T^2 = 36.0/346^2 = 3.00 \times 10^{-4}$ pada suhu rata-rata
geometris $\sqrt{300 \times 400} = \qty{346}{K}$.
\end{solution}

\begin{solution}{exo:b2:reaction-free-energy:8}
$\Delta_r S^\circ = 186.25 - 5.74 - 2(130.68) = \qty{-80.8}{J/(K.mol)}$,
$\Delta_r G^\circ = -74.87 + 298.15 \times 0.0808 = \qty{-50.8}{kJ/mol}$,
yaitu $\Delta_f G^\circ$ metana yang ditabelkan. Karena kedua suku bertanda
sama, metana menjadi tidak stabil di atas $T_i = 74\,870/80.8 \approx
\qty{926}{K}$ (\omterm{def:b2:reaction-free-energy:ellingham-approximation}{pendekatan Ellingham}): metana panas terengkah menjadi karbon
dan hidrogen, sebuah jalur untuk memperoleh keduanya.
\end{solution}

\begin{solution}{exo:b2:reaction-free-energy:9}
Substitusi: \ce{CH3CHBrCH3 + OH- -> CH3CH(OH)CH3 + Br-}; eliminasi:
\ce{CH3CHBrCH3 + OH- -> CH2=CHCH3 + H2O + Br-}. Selisih kedua \omterm{def:b2:reaction-free-energy:reaction-gibbs}{energi Gibbs
reaksi standar} itu adalah $\Delta(\Delta_r G^\circ) = 40 - T \times
0.090$, negatif di atas $40\,000/90 \approx \qty{440}{K}$. Eliminasi
menghasilkan satu molekul lebih banyak daripada substitusi: entropinya
lebih besar, dan suku entropi bertambah sebanding dengan $T$.
\end{solution}

\begin{solution}{exo:b2:reaction-free-energy:10}
Masing-masing dari $N_A$ molekul mempunyai 2 orientasi: $W = 2^{N_A}$,
sehingga $S =
k_B\ln 2^{N_A} = R\ln 2 = \qty{5.76}{J/(K.mol)}$. Kristal
itu tidak tersusun sempurna: pada suhu rendah molekul-molekulnya membeku
dalam orientasi acaknya dan tidak dapat mencapai satu-satunya susunan
teratur yang diandaikan oleh hukum ketiga.
\end{solution}

\begin{solution}{exo:b2:reaction-free-energy:11}
Dengan mengintegralkan $\dd\Delta_r S^\circ/\dd T = c/T$, diperoleh $\Delta_r S^\circ(T) = \Delta_r
S^\circ(T_0) + c\ln(T/T_0)$; Kirchhoff memberikan $\Delta_r H^\circ(T) = \Delta_r
H^\circ(T_0) + c(T - T_0)$; kurangkan $T\Delta_r S^\circ(T)$. Untuk batu
kapur pada \qty{1100}{K}, \omterm{def:b2:reaction-free-energy:ellingham-approximation}{pendekatan Ellingham} memberikan $178.32 - 1100 \times 0.16064 = \qty{1.62}{kJ/mol}$; koreksinya adalah $c[T - T_0 - T\ln(T/T_0)] = -0.00195 \times (801.85 - 1436.3) =
\qty{+1.24}{kJ/mol}$: $\Delta_r G^\circ = \qty{2.86}{kJ/mol}$. Galatnya, \qty{1.2}{kJ/mol}, kurang dari satu persen
$\Delta_r H^\circ$, tetapi di dekat \omterm{def:b2:reaction-free-energy:inversion-temperature}{suhu inversi} galat itu menentukan
tandanya.
\end{solution}

\begin{solution}{exo:b2:reaction-free-energy:12}
(a) $\Delta_r H^\circ = -763.2 + 944.0 = \qty{+180.8}{kJ/mol}$, $\Delta_r
S^\circ = 353.2 + 205.15 - 50.62 - 2(223.08) = \qty{+61.6}{J/(K.mol)}$;
$\Delta_r G^\circ(1000) = 180.8 - 61.6 = \qty{+119.2}{kJ/mol}$: \omterm{def:b2:reaction-free-energy:exergonic}{endergonik},
tidak ada klorinasi. (b) Untuk \ce{C + O2 -> CO2}, $\Delta_r G^\circ(1000) =
-393.5 - 2.9 = \qty{-396.4}{kJ/mol}$. Jumlahnya, \ce{TiO2(s) + 2Cl2(g) + C(s)
-> TiCl4(g) + CO2(g)}, mempunyai $\Delta_r G^\circ = \qty{-277.2}{kJ/mol}$:
karbon mengonsumsi dioksigen yang dihasilkan reaksi pertama, dan reaksi
gabungannya sangat \omterm{def:b2:reaction-free-energy:exergonic}{eksergonik}.
\end{solution}

\begin{solution}{pb:b2:reaction-free-energy:1}
\textbf{1.} \ce{CaCO3(s) -> CaO(s) + CO2(g)}.
\textbf{2.} $\Delta_r H^\circ = -635.09 - 393.51 + 1206.92 =
\qty{+178.3}{kJ/mol}$: \omterm{def:b2:reaction-enthalpy:exothermic}{endoterm}.
\textbf{3.} $\Delta_r S^\circ = 39.75 + 213.79 - 92.9 = \qty{+160.6}{J/(K.mol)}$,
positif karena gas terbentuk dari padatan.
\textbf{4.} $\Delta_r G^\circ = 178.32 - 298.15 \times 0.16064 =
\qty{+130.4}{kJ/mol}$.
\textbf{5.} Ya: penguraian memerlukan $\Delta_r G < 0$, dan di sini nilainya
sangat positif; karbon dioksida di udara, yang jauh di bawah \qty{1}{bar},
menurunkan $\Delta_r G$ tetapi jauh kurang dari \qty{130}{kJ/mol} (bab
berikutnya).
\textbf{6.} $\Delta_r C^\circ_p = 42.80 + 37.13 - 81.88 =
\qty{-1.95}{J/(K.mol)}$, sangat kecil: $\Delta_r H^\circ$ dan
$\Delta_r S^\circ$ hampir tidak berubah terhadap $T$.
\textbf{7.} $\Delta_r G^\circ(T) = 178.32 - 0.16064\,T$ (\unit{kJ/mol}).
\textbf{8.} \qty{+17.7}{kJ/mol} pada \qty{1000}{K}; \qty{-30.5}{kJ/mol}
pada \qty{1300}{K}.
\textbf{9.} $T_i = 178\,320/160.64 = \qty{1110}{K}$.
\textbf{10.} $\Delta_r G^\circ/T = 178.32/T - 0.16064$, yang turunannya
$-178.32/T^2 = -\Delta_r H^\circ/T^2$.
\textbf{11.} Pada $T_i$, $c[T_i - T_0 - T_i\ln(T_i/T_0)] = -0.00195 \times
(811.9 - 1459.4) = \qty{+1.26}{kJ/mol}$; $\Delta_r G^\circ$ mencapai nol
$1260/160.6 \approx \qty{8}{K}$ lebih tinggi, pada sekitar
\qty{1118}{K}: kurang dari satu persen.
\textbf{12.} Pada \qty{1000}{K}, $\Delta_r G > 0$: penguraian tidak dapat
berlangsung (kapur tohor justru akan menyerap karbon dioksida). Pada
\qty{1300}{K}, $\Delta_r G <
0$: batu kapur terurai.
\textbf{13.} $\Delta_r G$ menjadi positif: kapur tohor mengikat karbon
dioksida dan kembali menjadi karbonat.
\textbf{14.} $\delta S_{\text{c}} = -\Delta_r G\,\dd\xi/T = 30\,506/1300 =
\qty{23.5}{J/K}$ per mol.
\textbf{15.} Ya: \cref{ch:b2:chemical-potential} akan menunjukkan bahwa
$\Delta_r G = \Delta_r G^\circ + RT\ln(p_{\ce{CO2}}/p^\circ)$, yang
menurun jika tekanan karbon dioksida di bawah \qty{1}{bar}, sehingga
penguraian sudah dimulai di bawah $T_i$.
\textbf{16.} $n = 10^6/56.08 = \qty{1.783e4}{mol}$.
\textbf{17.} $1.783 \times 10^4 \times 178.32 = \qty{3.18e6}{kJ}$, yaitu
\qty{3.18}{GJ} per ton.
\textbf{18.} $1.783 \times 10^4 \times 44.01 = \qty{785}{kg}$ \ce{CO2}.
\textbf{19.} Kalor yang harus disediakan $3.18 \times 10^6/0.50 = \qty{6.36e6}{kJ}$, sehingga $6.36 \times 10^6/802.3 = \qty{7.93e3}{mol}$
metana, \qty{127}{kg}.
\textbf{20.} $7.93 \times 10^3 \times 44.01 = \qty{349}{kg}$ \ce{CO2}.
\textbf{21.} $785 + 349 = \qty{1134}{kg}$ per ton kapur tohor, 69\,\% di
antaranya dari batu kapur.
\textbf{22.} Karbon dioksida dari batu kapur ditentukan oleh
stoikiometri, apa pun yang memanaskan tungku; hanya bagian dari bahan
bakar yang dapat dikurangi.
\textbf{23.} Batu kapur terurai di bawah \qty{1}{bar} karbon dioksida di
atas \omterm{def:b2:reaction-free-energy:inversion-temperature}{suhu inversinya}, $\boldsymbol{T_i \approx \qty{1.11e3}{K}}$ (sekitar
\qty{840}{\celsius}).
\end{solution}
